\section{The self-referential network model}
\label{sec:model}

We work with a linearized model in which a sparse lunar constellation---orbiters and
surface nodes from more than one provider---estimates its own state from the ranges and
clock ties its nodes form with one another, with \emph{no} node position or time known
externally. This is the free (self-referential) network, the setting P4 identified as the
carrier of the coupled datum$\oplus$timescale gauge \citep{baweja2026gauge}; here it is the
setting in which an autonomous integrity monitor must work. This section fixes the state,
the observation model, the two subspaces that the rest of the paper keeps rigidly apart, the
fault taxonomy, and the kinematic scope. All quantities named \code{like\_this} are the
corresponding items in the open engine module \code{lunar\_faultobs}
(Section~\ref{sec:availability}).

\subsection{Nodes, state, and the per-node coordinates}
\label{sec:state}

The network has $M=8$ nodes: five near-side surface nodes (surface transponder beacons) and
three relay orbiters, drawn from at least two providers. Each node $j$ carries a local
five-parameter deviation coordinate
\begin{equation}
  \xvec_j \;=\; \bigl[\,\underbrace{\delta\bm p_j}_{\in\,\R^3},\;
                       \underbrace{\dtau_j}_{\text{offset}},\;
                       \underbrace{\dalpha_j}_{\text{rate}}\,\bigr]\;\in\;\R^{5},
  \qquad
  \xvec \;=\; [\,\xvec_1;\dots;\xvec_M\,]\;\in\;\R^{5M}=\R^{40},
  \label{eq:state}
\end{equation}
a position deviation $\delta\bm p_j$ (metres) and a clock offset $\dtau_j$ and rate
$\dalpha_j$. With $M=8$ the state dimension is $5M=40$; in the engine the layout constants
are \code{PARAMS\_PER\_NODE}${}=5$, \code{state\_dim}${}=40$.

\subsection{The differential observation model}
\label{sec:obsmodel}

Observations are inter-node \emph{differential} one-way pseudoranges and clock ties. A
differential row is built so that a whole-network Helmert transform (translation, rotation)
\emph{and} a common clock offset and rate cancel identically---both endpoints of an
inter-node measurement move together---so the row is invariant under the network's own datum
freedom. The engine builders \code{differential\_range\_row} and
\code{differential\_clock\_tie\_row} form these rows (the free-network analogue of P4's
anchored \code{oneway\_range\_row}); the spatial partials are computed through the
\emph{real} DE440 lunar physical-axis orientation path, so the line-of-sight geometry is the
external ingredient, while the node states are \emph{estimated}, not pinned.

Each of the $\binom{M}{2}=\binom{8}{2}=28$ node pairs is observed at three clock-integration
windows $T\in\{21600,43200,86400\}$~s (six hours, twelve hours, one day), giving
\begin{equation}
  n \;=\; \binom{8}{2}\times 3 \;=\; 28\times 3 \;=\; 84
  \label{eq:nrows}
\end{equation}
rows. Stacking the rows gives the linearized model
\begin{equation}
  \bm z \;=\; \Gmat\,\xvec \;+\; \bvec \;+\; \bm\varepsilon,
  \qquad
  \bm\varepsilon\sim\mathcal N(\bm 0,\Sig),
  \qquad
  \Wmat \;:=\; \Sig^{-1}\;=\;\operatorname{diag}(1/\sigma_i^2),
  \label{eq:model}
\end{equation}
with $\Gmat\in\R^{84\times 40}$ the network Jacobian, $\bvec\in\R^{84}$ a fault (bias)
vector, and $\Wmat$ the weight matrix. The measurement standard deviations cycle through
$\sigma\in\{0.5,1,2\}$~m, so $\Wmat\neq\bm{I}$ is a genuinely non-uniform weighting---which,
as Section~\ref{sec:obs} shows, is exactly the regime in which the correct minimum
detectable bias form matters. The epoch is 2024-01-01 TT (\code{t\_tt\_jc}${}=0.239986$
Julian centuries), fixing the DE440 orientation.

\subsection{The gauge: what the self-referential network cannot observe}
\label{sec:gauge}

Because no node is externally pinned, the network Jacobian is rank-deficient. Its null space
is the datum$\oplus$timescale \emph{gauge} of P4 \citep{baweja2026gauge}, realized
differentially:
\begin{equation}
  \Nsp(\Gmat)
  \;=\;\operatorname{span}\{\,
    \underbrace{3\text{ translations}}_{\text{rigid}},\;
    \underbrace{3\text{ rotations}}_{\text{rigid}},\;
    \underbrace{\text{common }\dtau}_{\text{offset}},\;
    \underbrace{\text{common }\dalpha}_{\text{rate}}\,\},
  \qquad \dim\Nsp(\Gmat)=8 .
  \label{eq:gauge}
\end{equation}
A common translation and infinitesimal rotation of the whole point set, a common clock
offset, and a common clock rate leave every internal differential observation unchanged, so
the eight corresponding directions are annihilated by $\Gmat$. Scale is \emph{observable},
not a gauge direction, because under the SI second and the $c$-defined metre a differential
range carries absolute length (the P4 result, imported unchanged). Since $\Wmat\succ0$,
$\Nsp(\Gmat^\top\Wmat\Gmat)=\Nsp(\Gmat)$, so the weighted normal matrix has the same
eight-dimensional gauge. With $\dim\Nsp(\Gmat)=8$ and state dimension $40$, the rank is
$\rank\Gmat=40-8=32$.

\subsection{Two subspaces, kept distinct}
\label{sec:twosubspaces}

The whole paper turns on holding two subspaces apart; we name them once, here.
\begin{itemize}
  \item The \textbf{parameter-space gauge}
  $\Nsp(\Gmat)\subset\R^{40}$ of \eqref{eq:gauge}: the unobservable \emph{state} moves---the
  absolute frame (P1) and the common clock offset and rate (P4). It lives in the state space
  $\R^{5M}$.
  \item The \textbf{measurement-space blind fault subspace}
  $\rng(\Gmat)\subset\R^{84}$: the \emph{biases} that leave the parity residual zero and are
  therefore undetectable (Section~\ref{sec:obs}). The common-mode error class of P3
  \citep{baweja2026commonmode} is a subset of this. It lives in the measurement space
  $\R^{n}$.
\end{itemize}
These are different subspaces of different spaces, with different dimensions ($8$ versus
$32$). What links them is that they are the two faces of the \emph{same} rank deficiency of
$\Gmat$, and Section~\ref{sec:obs} makes the link precise: an undetectable fault in
$\rng(\Gmat)$ is absorbed into a state error defined only modulo $\Nsp(\Gmat)$. We never
write that the two subspaces are ``the same.''

\subsection{Fault taxonomy}
\label{sec:taxonomy}

A fault is a bias vector $\bvec\in\R^{84}$; its \emph{support} distinguishes three classes,
of increasing adversarial power.
\begin{description}
  \item[Measurement fault.] A bias on a single measurement row: $\bvec=\gamma\,\bm e_i$ for
  one row $i$. This is the classical RAIM single-fault model.
  \item[Peer (Byzantine) fault.] A node $j$ corrupts, arbitrarily and possibly in collusion
  with other nodes, \emph{all} the measurements incident to it. Its signature is the
  column-block $\Bmat_j\in\R^{84\times m_j}$ selecting node $j$'s $m_j$ incident rows; the
  injected bias is $\bvec=\Bmat_j\bm v$ for arbitrary $\bm v\in\R^{m_j}$. In our network a
  node participates in $7$ pairs $\times\,3$ windows $=21$ rows, so a broad-coverage peer
  owning all of them has $m_j=21$. This is the Byzantine model of \citep{lamport1982},
  used for the fault \emph{model} only.
  \item[Provider fault.] A provider-wide ephemeris or timescale error is a \emph{coordinated}
  bias across all of that provider's nodes---a structured, multi-node fault treated in
  Section~\ref{sec:holdover}.
\end{description}

\subsection{Kinematic (snapshot) scope}
\label{sec:scope}

$\Gmat$ is the \emph{instantaneous geometry} Jacobian at the epoch: the gauge \eqref{eq:gauge}
of eight is the kinematic / snapshot datum defect. This is the scope of the paper, and it is
a real limitation. LiAISON \citep{hillborn2007} shows that \emph{orbital dynamics} over an
arc can recover absolute state from crosslinks by exploiting a dynamically-unique orbit; such
dynamical information can partially collapse $\Nsp(\Gmat)$. Our results are the geometry /
snapshot regime; the dynamical regime is complementary and is not claimed here. The three
integration windows in \eqref{eq:nrows} carry \emph{relative} clock information (per-node
rates become observable across windows), but not the \emph{common} rate, which remains a
gauge direction (Section~\ref{sec:holdover})---consistent with the kinematic scope.
